\subsection{معماری و پیاده‌سازی}
موتور فشرده‌سازی \lr{MoDeGPT} در فایل \lr{Codes/framework/compression\_engine/modegpt\_compressor.py} تعبیه شده و شامل مؤلفه‌های زیر است:
\begin{itemize}
    \item \textbf{لایه خطی \lr{MoDeGPTLoRALinear}:} لایه جایگزین که وزن‌های پایه فاکتورگیری‌شده را در فرمت استاندارد نگهداری کرده و با لایه‌های \lr{Conv1D} در مدل‌های سبکی همچون \lr{GPT-2} سازگاری کامل دارد. همچنین در صورت نیاز به فاین‌تیونینگ، قابلیت اضافه کردن ادپترهای \lr{LoRA} فعال می‌گردد.
    \item \textbf{کلاس \lr{MoDeGPTCompressor}:} مدیریت خوشه‌بندی لایه‌ها به جفت‌های مکمل (\lr{Joint Pairs})، بررسی سازگاری ابعاد لایه‌ها، مدیریت لایه‌های منفرد که جفت مناسبی ندارند توسط \lr{\_decompose\_single}، محافظت از پروجکشن‌های حساس در معماری‌های با مکانیزم توجه نامتقارن (\lr{GQA-aware skipping}) و امکان بازگشت امن به وزن اصلی در صورت نقض آستانه خطا.
\end{itemize}

\subsection{محدودیت‌ها و تحلیل شکست}
تحلیل فنی فرآیند فشرده‌سازی با \lr{modegpt\_compressor} نشان‌دهنده نکات و چالش‌های ساختاری زیر است:
\begin{itemize}
    \item \textbf{پیچیدگی جفت‌سازی در معماری‌های ناهمگن:} پیش‌فرض الگوریتم \lr{MoDeGPT} وجود جفت‌های متقارن خطی با ابعاد ورودی یکسان است. در مدل‌های مدرنی که از \lr{GQA} استفاده می‌کنند، ابعاد پروجکشن کلید و مقدار بسیار کوچکتر از پروجکشن پرسش است که امکان الحاق مستقیم بدون پدینگ یا نگاشت‌های اضافی را از بین برده و منجر به انصراف موتور از فشرده‌سازی جفت می‌گردد.
    \item \textbf{تثبیت رتبه سراسری بدون آگاهی از حساسیت طیفی:} تنظیم یک رتبه ثابت (\lr{rank=64}) برای کلیه زوج لایه‌ها بدون در نظر گرفتن محتوای اطلاعاتی طیف مقادیر منفرد می‌تواند در لایه‌های میانی که دربرگیرنده بیشترین پیچیدگی معنایی هستند منجر به افت دقت گردد.
\end{itemize}
